\subsection{应用于线性方程}

考虑交换环 A 上的一个含 \(n\) 个未知元的 \(n\) 个标量线性方程的方程组（II，§ 2，第 8 号）：

\[
\sum_ {j = 1} ^ {n} \lambda_ {i j} \xi_ {j} = \eta_ {i} \quad (1 \leqslant i \leqslant n).\tag{34}
\]

设 \(L\) 是 \((\lambda_{ij})\) 的 \(n\) 阶方阵；按通常方式把由 \(\xi_i\) （相应地， \(\eta_i\) ）组成的一列矩阵与元素 \(x = (\xi_i)\) （属于 \(A^n\) ）（相应地，元素 \(y = (\eta_i)\) （属于 \(A^n\) ））等同起来，方程组 (34) 也可写成（II，§ 10，第 3 号，命题 2）

\[
L. x = y.\tag{35}
\]

设 \(u\) 是自同态 \(x \mapsto L.x\) （ \(A^n\) 的，以 \(L\) 为关于典范基的矩阵）；说方程 (35) 对所有 \(y \in A^n\) 都有（至少）一个解，就是说 \(u\) 是满射；于是第 2 号的定理 1 蕴含以下命题：

命题13. 对交换环上的一个含 n 个未知元的 n 个线性方程的方程组，无论右端如何取值都至少有一个解的必要且充分条件是：方程组的矩阵的行列式可逆；此时方程组有唯一解。

若det L在A中不是零因子，则方程(34)等价于方程

\[
(\det L) L. x = (\det L) y.
\]

若 \(M\) 是 \(L\) 的代数余子式矩阵，则我们由 (34) 及第 6 号公式 (26) 导出关系式

\[
(\det L) x = ^ {t} M. y\tag{36}
\]

也可以写作

\[
(\det L) \xi_ {i} = \sum_ {j = 1} ^ {n} (- 1) ^ {i + j} (\det L ^ {i j}) \eta_ {j} = \det L _ {i} \quad (1 \leqslant i \leqslant n)\tag{37}
\]

其中 \(L_{i}\) 表示把 L 的指标为 i 的列换成 y 后所得的矩阵。公式 (37) 称为方程组 (34) 的克拉默公式；(34) 的每个解也是 (37) 的解。反之，由 (36) 并考虑到第 6 号公式 (28)，我们得出

\[
(\det L) (L. x - y) = 0\tag{38}
\]

因此，若det L在A中不是零因子，则方程组(34)与(37)等价；若det L可逆，则(34)的唯一解由下式给出

\[
\xi_ {i} = (\det L) ^ {- 1} (\det L _ {i}) \quad (1 \leqslant i \leqslant n).\tag{39}
\]

满足 det L 可逆的方程组 (34) 也称为克拉默方程组。

特别地，设 \(y = 0\) ；则由第 2 号命题 3 可得：

命题14. 对交换环上的一个含 n 个未知元的 n 个方程的齐次线性方程组，存在非零解的必要且充分条件是其矩阵的行列式为零因子。

